\section{基于线性关系的分类}
\label{sec:classification-geometry}

线性得分用特征的加权和比较类别。本节考察三种确定权重的依据：LDA使用类均值与协方差，感知机根据分类错误更新，SVM通过间隔与损失选择边界。核方法随后将内积换成特征空间内积，使相同的间隔准则能够表达非线性关系；线性表示的基础见第\ref{sec:classic-overview-linear}节。

LDA兼有投影与概率解释：这里推导共享协方差形式，第\ref{sec:classification-probability}节再以GDA模型族比较不同协方差设定。Fisher准则评价投影后的类统计量，感知机的有限更新保证依赖正间隔可分性，软间隔SVM则在目标中允许间隔违反。三者的条件分别服务于投影、更新和优化，比较时需要保留各自的对象。

\input{tex/04-classic-models/classification/01-lda-perceptron}
\input{tex/04-classic-models/classification/01-svm}

\subsection{模型对比与总结}
\label{sec:classification-geometry-summary}

三种基本模型都以线性得分比较类别，权重与阈值的来源却不同。生成式LDA估计类均值、协方差和先验，再用贝叶斯法则形成边界；Fisher判别直接最大化类间散度相对于类内散度的比例。感知机在违反分类条件时更新，可分训练集上的输出是随访问顺序变化的可行分隔面。SVM将间隔写进优化目标，以范数和训练损失的权衡选取解。表\ref{tab:classification-linear-comparison}比较这些依据，后面的扩展分别调节协方差、表示、稀疏化与间隔参数。

\begin{table}[htbp]
  \centering
  \small
  \begin{tabularx}{\linewidth}{@{}>{\RaggedRight\arraybackslash}p{18mm}>{\RaggedRight\arraybackslash}X>{\RaggedRight\arraybackslash}X@{}}
    \toprule
    模型 & 边界由什么决定 & 主要条件与代价 \\
    \midrule
    LDA & 类均值、共享协方差与类先验 & 高斯假设支持概率解释；协方差估计需足够样本或正则化。 \\
    感知机 & 错分样本触发的权重更新 & 正间隔可分性支持有限更新界；不可分时需停止规则或变体。 \\
    线性 SVM & 范数与间隔损失的权衡 & 硬间隔要求可分；软间隔需选择惩罚强度。 \\
    核 SVM & 特征空间中的间隔 & 核定义相似性；训练与预测开销受样本数和支持向量数影响。 \\
    \bottomrule
  \end{tabularx}
  \caption{线性与核化分类器的比较。LDA 一行指共享协方差的高斯生成模型；Fisher 投影本身不要求高斯分布。各模型的理论保证分别针对估计、更新次数或风险，不能直接互换。}
  \label{tab:classification-linear-comparison}
\end{table}

如果不同类别的类内变化方向和尺度明显不同，共享协方差就会限制LDA的表达能力。第\ref{sec:classification-lda-model}节已经定义的二次判别分析（QDA）保留类条件高斯假设，却为每个类别估计独立的协方差矩阵$\hat{\bm{\Sigma}}_c$。删去各类共同的常数项后，其判别得分为
\begin{equation}
  \begin{aligned}
    \delta_c(\bm{x})
      ={}&\log\hat{\pi}_c-\frac12\log\det\hat{\bm{\Sigma}}_c\\
        &-\frac12(\bm{x}-\hat{\bm{\mu}}_c)^\top
          \hat{\bm{\Sigma}}_c^{-1}(\bm{x}-\hat{\bm{\mu}}_c).
  \end{aligned}
  \label{eq:classification-qda-score}
\end{equation}
任意两类得分之差含有关于$\bm{x}$的二次项，所以边界一般是二次曲面。当协方差相同时，二次项相消，又得到 LDA。增加自由度也增加了估计负担：共享协方差模型需要$O(d^2+Cd)$个参数，QDA 则需要$O(Cd^2)$个参数；每一类只有少量样本时，独立协方差容易不稳定甚至奇异。第\ref{sec:classification-gda}节将从生成分布的角度统一这两种模型。

在完全共享与完全独立之间，可以让各类借用一部分合并统计量。\term{正则化判别分析}（regularized discriminant analysis, RDA）通过收缩协方差来控制这种取舍。Friedman 的原始方法包含两个调节方向：将类内协方差向合并协方差收缩，再向球形协方差收缩\scite{friedman1989rda}。
记$n_c$为第$c$类样本数，$\hat{\bm{\Sigma}}$为以$n$为分母的合并协方差。一种与原始构造相符的写法是
\begin{equation}
  \begin{aligned}
    \bm{A}_c(a)
      &=\frac{(1-a)n_c\hat{\bm{\Sigma}}_c+an\hat{\bm{\Sigma}}}
              {(1-a)n_c+an},\\
    \tilde{\bm{\Sigma}}_c(a,b)
      &=(1-b)\bm{A}_c(a)
        +b\,\frac{\operatorname{tr}\bm{A}_c(a)}{d}\bm{I},
        \qquad a,b\in[0,1].
  \end{aligned}
  \label{eq:classification-rda}
\end{equation}
当$b=0$时，$a=0$与$a=1$分别对应 QDA 与 LDA 的协方差估计；增大$b$会减弱各方向方差差异，并在平均方差为正时帮助获得可逆估计。两个参数应在训练集内部通过验证选择。当合并协方差和各类协方差均奇异时，可逆性还需要平均方差为正的球形收缩项或其他正则化。

Fisher 判别还可以沿核方法的路线扩展。Mika 等提出的\term{核 Fisher 判别}（kernel Fisher discriminant, KFD）在特征空间中重新计算类均值与散度，再最大化投影后的类间与类内散度比\scite{mika1999kernel}。
将投影方向写成$\bm{w}=\sum_{i=1}^n a_i\phi(\bm{x}_i)$，即可把目标表达为关于系数$\bm{a}$和核矩阵的广义 Rayleigh 商，并用正则化处理类内散度的奇异性。它与核 SVM 的区别仍在训练准则：KFD 关注整体类统计量，SVM 关注间隔及其违反程度。KFD 的基本解通常没有支持向量式的稀疏性，预测可能需要访问全部训练样本；若要从投影得分得到后验概率，还需额外建立得分分布或校准模型。

除了用间隔约束产生稀疏性，还可以为不同基函数分别学习先验的收缩强度。\term{相关向量机}（relevance vector machine, RVM）采用
\[
  g(\bm{x})=w_0+\sum_{i=1}^n w_i K(\bm{x}_i,\bm{x}),
  \qquad
  p(w_i\mid a_i)=\mathcal{N}(0,a_i^{-1}),
\]
其中$a_i>0$是每个权重各自的先验精度，即先验方差的倒数。分类时再用适当的概率链接函数把$g(\bm{x})$转换为类别概率。RVM通过最大化对权重积分后得到的边缘似然来估计各个精度，部分$a_i$可能趋向无穷，对应权重的分布随之集中到零。在数值实现中，可以按预定的精度或贡献准则剔除相应基函数，使预测只保留一部分训练样本的基函数。有限但很大的精度只造成强收缩，并不自动把权重变成精确的零；这里的稀疏化依赖精度学习及剔除过程。Tipping 系统介绍了这一稀疏贝叶斯学习框架\scite{tipping2001rvm}。
RVM输出模型内的预测分布，其调节对象是基函数及先验精度；精度学习包含非凸优化，分类似然还需要近似推断。相关向量的数量取决于任务与剔除准则，原模型的概率也需要在实际分布上检验校准。

如果希望调参时直接联系间隔错误与支持向量比例，可以考虑\term{$\nu$-SVM}（nu-support vector machine）。它把间隔阈值$\rho\geq0$也作为变量，在目标中用参数$\nu$奖励更大的阈值，同时惩罚平均松弛量：
\begin{equation}
  \begin{aligned}
    \min_{\bm{w},b,\rho,\bm{\xi}}\quad&
      \frac12\lVert\bm{w}\rVert_2^2-\nu\rho
      +\frac1n\sum_{i=1}^n\xi_i\\
    \text{满足}\quad&
      y_i(\bm{w}^{\top}\phi(\bm{x}_i)+b)\geq\rho-\xi_i,\\
    &\xi_i\geq0,\qquad \rho\geq0.
  \end{aligned}
  \label{eq:classification-nu-svm}
\end{equation}
在最优$\rho>0$时，其对偶乘子满足$0\leq\alpha_i\leq1/n$且$\sum_i\alpha_i=\nu$。因此至少需要$\nu n$个非零乘子，而每个严格违反间隔的样本都占用一个达到$1/n$的乘子，于是间隔错误比例不超过$\nu$，支持向量比例不低于$\nu$。这里计数的是训练集上的间隔错误，包括正确分类但落在间隔内的样本\scite{scholkopf2000nu}。

参数范围还受到类别比例约束。由偏置的驻点条件$\sum_i\alpha_i y_i=0$，正负两类的乘子和都为$\nu/2$；若两类样本数分别为$n_+$和$n_-$，盒约束便要求
\[
  0<\nu\leq\nu_{\max}
  :=\frac{2\min(n_+,n_-)}{n}.
\]
这里的上界不只是$\rho>0$解释所需的附加条件。若$\nu>\nu_{\max}$，令$\bm{w}=\bm{0}$、让常数偏置把全部样本判为多数类，并令少数类松弛量等于$2\rho$、多数类松弛量为零，就得到一条满足约束的射线；其目标值为$(\nu_{\max}-\nu)\rho$，随$\rho\to\infty$趋于$-\infty$。因此，虽然约束集合仍非空，优化问题却无界而没有有限最优解。在允许区间内取得有限最优解后，前述比例结论仍以最优$\rho>0$为条件；$\rho=0$时不能直接套用。

协方差收缩控制估计自由度，核化改变可表达的关系，RVM以先验精度及剔除实现稀疏化，$\nu$-SVM把训练间隔比例联系到调节参数。比较这些扩展时，固定特征处理和评价方式，再考察样本能否支持增加的自由度，以及训练、存储和预测成本是否符合用途。